Many problems arising in mathematics, computer science, biology and other disciplines can be modeled and analyzed through graph theory. A graph is a structure composed of vertices (or nodes) together with edges that link pairs of these vertices. Vertices correspond to entities or concepts, and edges capture the connections or interactions that exist between them. In this thesis, we work exclusively with undirected graphs, meaning every edge goes both ways and the connection it models is symmetric.

\par A graph $G$ is formally described by two sets: a vertex set $V = \{0, 1, \ldots, n - 1\}$, where $n$ is the number of vertices, and an edge set $E \subseteq V \times V$ whose elements are pairs of vertices joined by an edge.

\par One of the central notions we rely on is that of an induced path. Let $S \subseteq V$ be a subset of vertices. By retaining only the vertices in $S$ and those edges from $E$ whose both endpoints belong to $S$, we construct the induced subgraph $G[S] = (S, E')$ of the original graph $G = (V, E)$. If this subgraph forms a simple path, then we have found an induced path within the original graph. In other words, if the vertices we have chosen have no shortcut edges, the only edges being those of consecutive vertices, then the subgraph of these vertices will be a path.

The second procedure, \textbf{INDUCED-PATHS-FROM-VERTEX} (Algorithm~\ref{alg:InducedPathsFromVertex}), takes as parameters the graph \( G \), a source vertex \( s \), the longest induced path found so far \( P_\text{max} \), and the path currently being built \( P_\text{temp} \). In line 1, vertex \( s \) is appended to \( P_\text{temp} \). Line 2 computes the neighbor set \( N(s) \) of \( s \) in \( G \). If this set is non-empty, the loop in lines 4 to 9 iterates over each neighbor. For every vertex \( t \in N(s) \), lines 5 to 7 build a temporary subgraph \( G_\text{temp} \) by removing \( s \) and all its neighbors except \( t \). Because every neighbor of \( s \) other than \( t \) is deleted, no vertex adjacent to \( s \) can later appear on the path and violate the induced property. The search space also shrinks with each such deletion. Line 8 recurses on \( G_\text{temp} \) with \( t \) as the new source to extend the path by one vertex. If \( s \) has no neighbors (\( N(s) = \emptyset \)), the path cannot grow further. Lines 11 to 13 check whether \( P_\text{temp} \) is longer than \( P_\text{max} \) and, if so, update \( P_\text{max} \). Line 15 then removes \( s \) from \( P_\text{temp} \) before the recursion returns, so that subsequent branches do not carry over vertices from the current one.
